\documentclass[11pt]{article}
\usepackage[margin=1in]{geometry}
\usepackage{amsmath,amssymb,amsthm}
\usepackage{physics}
\usepackage{siunitx}
\usepackage{graphicx}
\usepackage{hyperref}

\title{System for Reaching Prime Numbers Trough Lines}
\author{Aaron Cuevas (Aaron-Cuevas)}
\date{2026-06-28}

\begin{document}
\maketitle
An insight looking trough prime numbers graphical arrangement
\section
Is possible to explain visually why some geometrical patterns appear in certain prime numer representations. Looking trough geometry it can be possible to explain the distribution of complex patterns. 
The course of ideas will follow until drawing over the known polar-spiral structure, lines, being the result of projecting residual classes.
\subsection
The prime numbers are commonly represented with the equations of the polar-spiral graph, where every prime number with radio *p*
and angle *p*:

\begin{equation}
    z_p = p e^{ip}
\end{equation}
And this behaviour come from a polar representation of integers, with a gap made between certain integers and the distribution of prime numbers in 
residual clases coprimes of module (N)
\subsection
Every natural number can be colocated in the plane trough
[
z_n = n e^{in}
]
This equals as saying
[
x_n = n\cos(n)
]
[
y_n = n\sin(n)
]
Where the number (n) determines simultaneously the distance to center and the rotation angle. 
Filtering only the prime numbers, we obtain:
[
z_p = p e^{ip}
]
Where (p) passes by the prime numbers.

\section{Partitions on the spiral}
For concreting the multiple-partitions spiral structure, we choose a module (N), reaching that every natural number 
can be calssified according to its residue.
[
n \equiv r \pmod N
]
And it's found that numbers of one residual calss have the form
[
n = r + kN
]
Putting then in an polar-spiral
[
z_{r+kN} = (r+kN)e^{i(r+kN)}
]
And separating the angular terms:
[
z_{r+kN} = e^{ir}(r+kN)e^{ikN}
]
Stabishing the relations with the nearby multiples, we have that if (N) is near a multiple of (2\pi), so:
[
N \approx 2\pi m
]
and by that:
[
e^{ikN} \approx 1
]
Implies that the numbers of one same residual class tends to be near the same angular direction.
By this, we have built the lines that emerge from center.
\subsection{Introduction of Dirichet Theorem}
The tendency of the residual classes for obtaining infinite prime numbers don't apply to all of them.
If a residue (r) shares a factor with (N), so a lot of numbers of the form
[
r+kN
]
Will be compound. For that are considered specially the classes where 
[
\gcd(r,N)=1
]
Being them the residual coprimal classes with (N). 
According with Dirichlet's theorem, if (r) and (N) and coprimes, so exist infinite coprimes 
in the aritmethic progression:
[
r, r+N, r+2N, r+3N,\dots
]
The lines corresponds to residual classes (r\pmod N), with
[
\gcd(r,N)=1
]
\subsection{Module 44}
Remembering the example given by the initial explanation where is inspired this method, 
we have that the number 44 is near of seven whole spins,
[
44 \approx 7(2\pi)
]
For that the residual classes of module (44) generates more numerous and defined lines.
The residues that can contain infinite primes are those coprimes with (44)
[
\gcd(r,44)=1
]
as:
[
\varphi(44)=20
]
There are twenty permited classes that corresponds to 20 lines expected in the visualization. 
\subsection{Error Tracking}
As the alineation ain't perfect because (N) doesn't equal exactly a multiple of (2\pi). It can be writed
[
N = 2\pi m + \varepsilon
]
Where (\varepsilon) is the angular error. Every time we step ahead one unity in the same residual class, the error acumulates:
[
k\varepsilon
]
If (\varepsilon) is little, the line seems straight during very long time. If its larger, the line bends faster.
So on, the straightnes of the lines depends on how good is the aproximation:
[
N \approx 2\pi m
]

/section{Proposed method}
Formulating the analysis method, we find that, in first, we graph whe prime numbers through the function:
[
z_p = p e^{ip}
]
Then, we choose a modulo (N) near a multiple of (2\pi):
[
    N\approx 2\pi m
]
And we divide the numbers in residual classes:
[
p \equiv r \pmod N
]
We identify the coprimal classes:
[
\gcd(r,N)=1
]
And trace geometrical guides to that classes, compare the visual distribution of the primes with the permited resudual classes, and measure the accumulated angular deviation:
[
\varepsilon = N - 2\pi m
]
We analice what modules produce more persistent lines.
% Domain: 🔬 Academic Systems
% Write your insight below. Examples you can delete:
%
% Inline math: the wavelength is $\lambda = \SI{632.8}{nm}$.
%
% \begin{equation}
%   E = \hbar\,\omega, \qquad \pdv{\psi}{t} = -\frac{i}{\hbar}\,\hat{H}\psi
% \end{equation}

\end{document}
